\chapter{Les maisons de market making d'options}\label{ch:in:the-options-market-making-houses}

Au début des années 1980, un groupe d'amis rencontrés à l'université commence à trader sur le parquet du Philadelphia Stock Exchange, et
en 1987, ils se réunissent au sein d'une même société. Deux ans plus tard, deux traders qui travaillent côte à côte sur le parquet
de la bourse d'options sur actions d'Amsterdam fondent une autre société. En 2002, une troisième société prend sa place dans la corbeille des options sur indices du parquet du
Chicago Board Options Exchange. Chaque maison raconte ainsi sa propre origine, sur ses propres pages, et chacune est aujourd'hui
un teneur de marché d'options qui cote sur écrans partout dans le monde. Le parquet leur a appris à évaluer chaque strike,
à tenir leur livre, et à couvrir ce qu'ils détenaient~; ils recrutent et forment encore des gens pour faire de même, désormais à la vitesse du
logiciel. Ce chapitre décrit ces maisons en tant qu'employeurs~: d'où elles viennent, comment elles forment leurs équipes, en quoi leur travail
se distingue de celui des sociétés électroniques du chapitre~2, et ce que montre celle d'entre elles qui publie ses résultats.

\section{Trois filiations~: Amsterdam, Chicago, Philadelphie}

Les options cotées ont été négociées autrefois sur les parquets des bourses, notamment à Chicago, Amsterdam et Philadelphie. Les teneurs de marché
de ces parquets étaient de petites associations de \omterm{def:in:proprietary-market-makers-and-high-frequency-firms:floor}{traders de parquet} (chapitre~2) qui évaluaient les options à l'aide
de modèles et les couvraient avec les actions sous-jacentes. Certaines d'entre elles sont devenues les maisons d'options d'aujourd'hui.

\begin{definition}[Maison de market making d'options]\label{def:in:the-options-market-making-houses:house}
Une \emph{maison de market making d'options}\index{maison de market making d'options} est une \omterm{def:in:a-map-of-the-industry:ptf}{firme de trading pour compte propre} (principal trading firm) dont l'activité principale
consiste à coter des options cotées en bourse, sur de nombreux strikes et de nombreuses échéances de nombreux sous-jacents, généralement au titre de
désignations comme teneur de marché assorties d'obligations de cotation, et à couvrir les positions qui en résultent dans le sous-jacent et dans d'autres options.
\end{definition}

La définition décrit une société~; le rôle qu'elle joue sur une plateforme de négociation, celui du teneur de marché d'options avec ses obligations de cotation
et ses protections, est traité dans le livre~1 (chapitre~24). La \Cref{fig:in:the-options-market-making-houses:lineage} montre les origines
que les maisons indiquent. Trois d'entre elles rattachent leur fondation à un parquet précis~; les autres donnent une ville et une année.

\begin{omfigure}
\begin{tikzpicture}[font=\footnotesize,>=stealth,
  venue/.style={draw,rounded corners,fill=omDef!12,minimum width=3.6cm,minimum height=0.6cm,align=center},
  firm/.style={draw,rounded corners,fill=omProp!8,minimum width=3.2cm,minimum height=0.5cm,align=center}]
\node[venue] (ams) at (0,0) {bourse d'options\\sur actions d'Amsterdam};
\node[venue] (phl) at (4.4,0) {Philadelphia\\Stock Exchange};
\node[venue] (cbo) at (8.8,0) {Chicago Board\\Options Exchange};
\node[firm] (opt) at (-1.0,-1.5) {Optiver (1986)};
\node[firm] (imc) at (-1.0,-2.3) {IMC (1989)};
\node[firm] (sig) at (4.4,-1.5) {Susquehanna (1987)};
\node[firm] (gro) at (4.4,-3.1) {Group One (1989)};
\node[font=\scriptsize,gray] at (4.4,-3.65) {San Francisco};
\node[firm] (pea) at (9.8,-1.5) {PEAK6 (1997)};
\node[firm] (bel) at (9.8,-2.3) {Belvedere (2002)};
\node[font=\scriptsize,gray] at (-1.0,-2.85) {Amsterdam};
\node[font=\scriptsize,gray] at (9.8,-2.85) {Chicago};
\draw[->,thick,omDef] (imc.east) -- ++(0.7,0) |- (ams.south east);
\draw[->,thick,omDef] (sig.north) -- (phl.south);
\draw[->,thick,omDef] (bel.west) -- ++(-0.6,0) |- (cbo.south west);
\end{tikzpicture}

\omcaption{Où six teneurs de marché d'options disent avoir commencé~: une flèche pleine signifie que la page propre de la société rattache sa fondation
à une négociation sur ce parquet~; les sociétés sans flèche n'indiquent qu'une ville et une année. Sources~: les pages des sociétés
(registre F1, F3--F5, F7, F8).}\label{fig:in:the-options-market-making-houses:lineage}
\end{omfigure}

Quatre des six sociétés donnent des dates de fondation dans les années 1980, la décennie des parquets~; une dans les années 1990 et
une dans les années 2000. Deux des six ont commencé à Amsterdam~; le chapitre~27 décrit les dispositifs fiscaux et de visa qui amènent le personnel vers
les villes où se trouvent les maisons.

\begin{definition}[Essaimage]\label{def:in:the-options-market-making-houses:spinout}
Un \emph{essaimage}\index{essaimage} (spin-out) est une société fondée par des personnes qui quittent une société existante pour exercer la même activité ou une activité connexe
à leur propre compte. Contrairement à une scission, la société mère ne distribue ni ne conserve les actions de la nouvelle société.
\end{definition}

Dans ce secteur, les filiations se racontent souvent comme des chaînes d'\omterm{def:in:the-options-market-making-houses:spinout}{essaimages}, d'une maison vers les traders qui l'ont quittée. Aucune des
six sociétés ne mentionne une telle origine sur ses propres pages, si bien que le graphe du chapitre ne comporte aucune arête d'\omterm{def:in:the-options-market-making-houses:spinout}{essaimage}~; une filiation bâtie à partir de sources
publiques est plus maigre que celle que le secteur raconte sur lui-même, et le chapitre~28 revient sur ce que le dossier public
montre réellement des personnes qui quittent leur maison pour créer des sociétés.

\section{Programmes de formation tels que les sociétés les décrivent}

Une maison d'options transforme des diplômés en traders qui évaluent et couvrent des milliers d'instruments à la fois, et elle le dit
lorsqu'elle recrute. Une maison, dont les fondateurs ont bâti leurs premières stratégies grâce à une «~expérience du poker~», affirme que «~nous utilisons des jeux de stratégie et la science de la décision pour
former beaucoup de nos nouveaux employés~»~; sa page de campus publie le témoignage d'un stagiaire sur un stage de dix semaines passé dans un desk de
volatilité d'indices. Une autre maison décrit son enseignement comme «~un apprentissage conçu pour les marchés en direct et les systèmes en temps réel~».
Le contenu enseigné est celui des livres~1, 5 et 11~: profils de gains des options et parité, grecques, surfaces de volatilité,
inventaire et couverture. Ce que les maisons ajoutent, c'est la pratique~: marchés simulés, jeux de trading et plusieurs mois de travail supervisé
avant qu'un trader ne détienne seul un risque.

\begin{remark}[Ce qu'est un programme de formation]\label{rem:in:the-options-market-making-houses:programme}
Un programme de formation est l'investissement de la société pour faire d'une recrue une personne à qui elle confie des risques. Sa durée, son
contenu et la question de savoir si une recrue qui part tôt doit rembourser quoi que ce soit relèvent des clauses du contrat de travail, et diffèrent selon
la société et le pays. Les pages des sociétés décrivent l'expérience, non les clauses~; un candidat lit le contrat
(chapitre~13).
\end{remark}

\section{En quoi une maison d'options se distingue d'une société delta-one électronique}

Les sociétés du chapitre~2 cotent surtout des instruments avec un prix chacun~: une action, un contrat à terme, une paire de devises. Une maison d'options
cote une surface. Les différences sont les suivantes.

\begin{enumerate}
\item \textbf{Le nombre d'instruments.} Un sous-jacent avec vingt échéances cotées et cinquante strikes chacune compte
$20\times50\times2=2\,000$ séries d'options. Une maison cotant quelques centaines de sous-jacents entretient des centaines de milliers de
cotations, mises à jour ensemble à mesure que le sous-jacent bouge (livre~11, chapitre~19).
\item \textbf{Le risque détenu.} Un teneur de marché delta-one vise à être neutre en quelques minutes. Une maison d'options détient un livre de
volatilité, de skew et de décote temporelle pendant des jours ou des semaines, couvrant la direction (delta) en continu et le reste
en négociant d'autres options.
\item \textbf{Les obligations.} Les désignations comme teneur de marché sur les bourses d'options emportent des obligations de cotation sur de nombreuses
séries (livre~1, chapitre~24)~: la maison doit être présente alors qu'elle préférerait ne pas l'être.
\item \textbf{Le modèle.} Les prix proviennent d'une surface de volatilité calibrée en temps réel sur le marché (livre~11, chapitre~19~; livre~5),
et non d'un prix juste unique~; le jugement d'un trader agit sur les paramètres de la surface.
\end{enumerate}

Pour un salarié, la conséquence est un autre métier. Un trader d'une maison d'options gère un livre de grecques et décide
où doit se situer la surface~; un trader d'une société delta-one supervise des stratégies qui se disputent un seul prix.
Les deux utilisent les mêmes ordinateurs~; les questions sont, elles, différentes.

\begin{example}[L'échelle d'une maison]\label{ex:in:the-options-market-making-houses:scale}
Une maison, fondée sur le parquet d'Amsterdam, affirme qu'elle négocie sur 120 plateformes avec 2\,070 salariés répartis dans dix bureaux.
Si chaque plateforme n'admet à la cote des options que sur une centaine de sous-jacents, chacun avec 2\,000 séries, la maison pourrait coter de l'ordre de
$120\times100\times2\,000=24$ millions de séries~; le nombre réel est plus petit, car de nombreuses plateformes admettent les mêmes
sous-jacents et de nombreuses séries ne sont jamais cotées~; mais l'ordre de grandeur explique pourquoi la cotation relève du logiciel et
le travail du trader relève de la supervision et du jugement de valorisation.
\end{example}

\section{Taille et empreinte}

Contrairement à la plupart des sociétés de trading privées, une maison d'options publie ses résultats chaque année, accompagnés d'une brève revue annuelle.

\begin{dated}{2026-09}{Les maisons d'options dans leurs propres termes}\label{dat:in:the-options-market-making-houses:houses}
\textbf{Optiver} (résultats 2025, publiés le 31 mars 2026)~: produit net de trading 4\,556~millions d'euros (2024~: 
3\,494~millions d'euros~; 2023~: 2\,773~millions d'euros)~; résultat net attribuable aux détenteurs de fonds propres 1\,769~millions d'euros (2024~: 
1\,369~millions d'euros~; 2023~: 1\,158~millions d'euros)~; fonds propres totaux 5\,490~millions d'euros (2024~: 4\,905~millions d'euros)~; «~plus de 2\,000
salariés~» (2\,100 dans le communiqué de 2024), fondée à Amsterdam en 1986, avec des bureaux dans dix autres villes.
\textbf{IMC}~: fondée en 1989~; 2\,070 salariés répartis dans dix bureaux, négociant sur 120 plateformes.
\textbf{Group One}~: fondée à San Francisco en 1989~; «~plus de 100~» personnes réparties dans quatre villes~; un teneur de marché d'options ou
spécialiste de chaque grand groupe de bourses américain. \textbf{Akuna Capital}~: «~une société de trading pour compte propre de premier plan, spécialisée
dans le market making d'options~». Les deux premières décrivent leur activité d'options dans leurs propres termes~: l'une valorise «~1M+~»
instruments «~des options et des ETF aux actions, obligations et devises~»~; l'autre se présente comme «~un des trois premiers
fournisseurs de liquidité en volume négocié sur options cotées dans le monde~».
\end{dated}

Les résultats publiés rendent visible l'économie d'une maison (la \cref{fig:in:the-options-market-making-houses:results}).
Le produit net de trading a progressé de 26~\% en 2024 et de 30~\% en 2025~; la marge nette sur ce produit a été de 39.2~\% et 38.8~\%~; la
rentabilité moyenne des fonds propres a été de 30.4~\% et 34.0~\%. Avec 2\,100 salariés en 2024, le produit net de trading atteignait 1.66~million d'euros
par salarié~; avec «~plus de 2\,000~» en 2025, au plus 2.28~million d'euros par salarié. Le chapitre~12 place ces chiffres à côté
des autres sociétés qui déposent leurs comptes.

\begin{omfigure}
\begin{tikzpicture}
\begin{axis}[width=10.5cm,height=5.2cm,ybar,area legend,bar width=14pt,xmin=2022.4,xmax=2025.6,ymin=0,ymax=6,
  xtick={2023,2024,2025},xticklabel style={/pgf/number format/1000 sep={}},ylabel={\small milliards d'euros},
  tick label style={font=\footnotesize},grid=major,grid style={gray!20},table/col sep=comma,
  legend style={legend cell align=left,font=\scriptsize,at={(0.5,-0.25)},anchor=north,/tikz/every even column/.append style={column sep=8pt}},legend columns=3]
\addplot[fill=omDef!55,draw=omDef] table[x=year,y=nti]{figdata/industry/03-the-options-market-making-houses/results.csv};
\addplot[fill=omThm!55,draw=omThm] table[x=year,y=profit]{figdata/industry/03-the-options-market-making-houses/results.csv};
\addplot[fill=gray!40,draw=gray] table[x=year,y=equity]{figdata/industry/03-the-options-market-making-houses/results.csv};
\legend{produit net de trading,résultat net,fonds propres totaux (fin d'exercice)}
\end{axis}
\end{tikzpicture}

\omcaption{Résultats publiés d'une maison d'options détenue en privé, 2023--2025~: le produit net de trading est passé de 
2.8~milliards d'euros à 4.6~milliards d'euros, et le résultat net de 1.2~milliard d'euros à 1.8~milliard d'euros. Données~: communiqués de résultats d'Optiver
pour 2024 et 2025, via \texttt{in\_options.results}.}\label{fig:in:the-options-market-making-houses:results}
\end{omfigure}

\begin{remark}[Pourquoi une maison publie]\label{rem:in:the-options-market-making-houses:why}
Une société détenue en privé qui publie ses résultats fait ce choix~: pour ses contreparties, ses prêteurs et ses recrues. Les
chiffres sont ceux de la société, audités mais non déposés dans un registre que le lecteur puisse vérifier ligne à ligne, et la revue est
rédigée par la société. Il s'agit d'une source primaire pour ce que la société déclare, et non pour sa comparaison avec les sociétés qui
ne publient rien.
\end{remark}

\section{Ce qu'indiquent les dépôts pour les maisons d'options}

Trois des teneurs de marché de l'échantillon, Optiver, IMC et Susquehanna, sont les maisons d'options de ce chapitre, et
ensemble ils déposent assez de demandes de certification de conditions de travail (chapitre~14) pour être comparés aux autres teneurs de marché sans
publier la rémunération d'une société en particulier.

\begin{dated}{2025-09}{Maisons d'options et autres teneurs de marché dans les dépôts}\label{dat:in:the-options-market-making-houses:lca}
Demandes de certification de conditions de travail de l'exercice 2025, salaire de base proposé~: les trois maisons d'options, 130 demandes, médiane 150\,000~\$
(10\up{e}--90\up{e}~centile 84\,800--200\,000~\$)~; les sept autres teneurs de marché, 237 demandes, médiane 200\,000~\$
(130\,000--300\,000~\$). Chercheurs 150\,000~\$ contre 200\,000~\$~; ingénieurs logiciels 146\,100~\$ contre 190\,000~\$~;
traders 200\,000~\$ contre 192\,500~\$. Au niveau de salaire~I~: 23 des 130 demandes des maisons d'options, 8 des 237 autres.
Salaire de base uniquement.
\end{dated}

\begin{omfigure}
\begin{tikzpicture}
\begin{axis}[width=10.5cm,height=6cm,xmin=50,xmax=320,ymin=0.4,ymax=8.6,ytick=data,
  yticklabels={{traders~; autres},{traders~; maisons d'options},{ingénieurs logiciel~; autres},{ingénieurs logiciel~; maisons d'options},{chercheurs~; autres},{chercheurs~; maisons d'options},{tous postes~; autres},{tous postes~; maisons d'options}},
  yticklabel style={font=\scriptsize},xticklabel style={font=\footnotesize},xlabel={\small salaire de base proposé, milliers de \$},
  grid=major,grid style={gray!20},table/col sep=comma]
\addplot[draw=none,mark=none,error bars/.cd,x dir=both,x explicit,error mark=none,error bar style={omDef,line width=1pt}]
  table[x=p50,y=pos,x error plus expr=\thisrow{p90}-\thisrow{p50},x error minus expr=\thisrow{p50}-\thisrow{p10}]
  {figdata/industry/03-the-options-market-making-houses/options_pay.csv};
\addplot[draw=none,mark=none,error bars/.cd,x dir=both,x explicit,error mark=none,error bar style={omDef!40,line width=6pt}]
  table[x=p50,y=pos,x error plus expr=\thisrow{p75}-\thisrow{p50},x error minus expr=\thisrow{p50}-\thisrow{p25}]
  {figdata/industry/03-the-options-market-making-houses/options_pay.csv};
\addplot[only marks,mark=|,mark size=4.5pt,line width=1.5pt,omThm] table[x=p50,y=pos]{figdata/industry/03-the-options-market-making-houses/options_pay.csv};
\end{axis}
\end{tikzpicture}

\omcaption{Salaire de base proposé dans les demandes de certification de l'exercice 2025~: les trois maisons d'options face aux autres
teneurs de marché, par famille de postes~: du 10\up{e} au 90\up{e}~centile (trait fin), intervalle interquartile (trait épais), médiane (repère). Regroupé sur
trois employeurs ou plus~; aucun chiffre d'une société n'est présenté. Données~: \texttt{data/industry/lca\_options.csv}, via
\texttt{in\_options.options\_pay}.}\label{fig:in:the-options-market-making-houses:lca}
\end{omfigure}

La médiane des salaires de base proposés par les maisons d'options est inférieure de 50\,000~\$ (intervalle bootstrap à 95~\%, entre $-\$53\,900$ et $-\$35\,000$) pour tous les intitulés
et pour les chercheurs, et celle des ingénieurs logiciels est inférieure de 43\,900~\$~; celle des traders est au même niveau que celle des autres
(la \cref{fig:in:the-options-market-making-houses:lca}). Deux précautions s'imposent avant toute lecture. Les maisons d'options déposent
beaucoup plus de demandes au niveau d'entrée, 17.7~\% au niveau de salaire~I contre 3.4~\%, ce qui est cohérent avec les programmes de formation destinés aux jeunes diplômés
décrits plus haut~; le salaire de base ne représente qu'une partie de la rémunération chez un teneur de marché, où le chapitre~14 a trouvé que le
rapport implicite entre rémunération totale et salaire de base est le plus élevé. Les dépôts montrent où se situent les salaires de base proposés, non ce que rapportent réellement les postes.

\section{Tutoriel~: les origines des maisons et les chiffres d'une maison}\label{tut:in:the-options-market-making-houses:tut}

\begin{tutorial}
\textbf{Objectif.} Construire le graphe de filiation à partir des déclarations des sociétés elles-mêmes, et calculer l'économie d'une maison à partir de ses
résultats publiés. \textbf{État final~:} les arêtes de la \cref{fig:in:the-options-market-making-houses:lineage} et les ratios de
la section ci-dessus.
\begin{tutsteps}
\item \textbf{Le graphe.} \texttt{data/industry/lineage.csv} recense trois plateformes de négociation et six sociétés~; trois lignes portent une arête
allant d'une société au parquet sur lequel ses fondateurs négociaient. \texttt{in\_options.graph()} le charge avec \texttt{firm.lineage},
qui refuse un nœud ou une arête sans ligne au registre
(le \cref{lst:in:the-options-market-making-houses:add}).
\omcode{code/firm/lineage/firm_lineage.py}{51}{61}{Une arête n'entre dans le graphe qu'avec une source et entre des
nœuds connus.}\label{lst:in:the-options-market-making-houses:add}
\item \textbf{Requête.} \texttt{g.by\_place()} regroupe les sociétés par ville, \texttt{g.by\_decade()} compte les décennies de fondation
(années 1980~: 4, années 1990~: 1, années 2000~: 1), et \texttt{g.unsourced()} recense les trois sociétés dont l'origine retenue dans la source publique
ne se rattache à aucun parquet.
\item \textbf{Les chiffres.} \texttt{in\_options.metrics()} calcule la croissance, la marge, la rentabilité moyenne des fonds propres et le produit
par salarié à partir de \texttt{data/industry/optiver\_results.csv}, en traitant «~plus de 2\,000 salariés~» comme une borne inférieure des
effectifs, et donc une borne supérieure du produit par salarié.
\end{tutsteps}
\end{tutorial}

\textbf{Pour aller plus loin.} Ajouter l'année de fondation et l'origine d'une autre maison d'options à partir de ses propres pages, et vérifier
si cela modifie le décompte par décennie~; calculer quel serait le produit par salarié de 2025 si l'effectif était de 2\,500
(exercice~7).

\section{Construction~: le graphe de filiation}\label{bld:in:the-options-market-making-houses:lineage}

\begin{build}
\bfield{Rôle} Une généalogie documentée des sociétés de trading, pour les parquets de ce chapitre et pour la question du chapitre~28 de savoir
d'où viennent les nouvelles sociétés.
\bfield{Interface} \texttt{firm.lineage}~: \texttt{Node(name, kind, year, place, ledger)},
\texttt{Edge(src, dst, relation, ledger)}~; \texttt{Graph.load(path)}, \texttt{add\_node}, \texttt{add\_edge}~; \texttt{firms}, \texttt{by\_place},
\texttt{by\_decade}, \texttt{children}, \texttt{ancestors}, \texttt{unsourced}.
\bfield{Règles} Aucun nœud ni aucune arête sans ligne de registre~; une arête relie deux nœuds connus~; une année de fondation est facultative et
n'est jamais déduite.
\bfield{Tests d'acceptation} \texttt{code/firm/lineage/tests/}~: regroupement par lieu et par décennie, enfants et ancêtres sur un
graphe construit~; les arêtes sans source et les arêtes pendantes sont refusées.
\bfield{Pour aller plus loin} Exporter le graphe en TikZ avec des positions calculées~; ajouter des arêtes d'acquisition à partir des dépôts réglementaires (le
formulaire 10-K d'une société cotée mentionne ses acquisitions).
\end{build}

\begin{omsources}
\item Les pages propres aux sociétés~: IMC (About us), Susquehanna (About, Campus), Belvedere Trading, PEAK6, Group One, Akuna
Capital, Optiver (page d'accueil, Education).
\item Optiver, communiqués de résultats pour 2024 (27 mars 2025) et 2025 (31 mars 2026).
\item Cboe Global Markets, Form 10-K for 2025, sur le parquet de négociation de Cboe Options.
\end{omsources}

\section{Exercices}

\begin{exercise}[$\star$]\label{exo:in:the-options-market-making-houses:1}
Combien de séries d'options un sous-jacent ayant 15 échéances et 40 strikes par échéance compte-t-il, en comptant les calls et les puts~?
\end{exercise}

\begin{exercise}[$\star$]\label{exo:in:the-options-market-making-houses:2}
À partir de l'encadré daté, calculer la marge nette sur le produit net de trading de la maison en 2023.
\end{exercise}

\begin{exercise}[$\star$]\label{exo:in:the-options-market-making-houses:3}
Parmi les six sociétés de la \cref{fig:in:the-options-market-making-houses:lineage}, lesquelles rattachent leur fondation à un parquet, et
à quel parquet~?
\end{exercise}

\begin{exercise}[$\star\star$]\label{exo:in:the-options-market-making-houses:4}
Calculer la croissance du résultat net de la maison en 2024 et en 2025, et expliquer pourquoi le résultat net a progressé plus lentement que le produit
net de trading en 2024.
\end{exercise}

\begin{exercise}[$\star\star$]\label{exo:in:the-options-market-making-houses:5}
Le produit net de trading par salarié de la maison s'élevait à 1.66~million d'euros en 2024. Dans quelle fourchette se situe le produit par salarié de 2025 si
l'effectif était compris entre 2\,000 et 2\,400~?
\end{exercise}

\begin{exercise}[$\star\star$]\label{exo:in:the-options-market-making-houses:6}
Donner deux raisons pour lesquelles une maison d'options garde le risque plus longtemps qu'un teneur de marché électronique delta-one.
\end{exercise}

\begin{exercise}[$\star\star\star$]\label{exo:in:the-options-market-making-houses:7}
\emph{Programmation.} Calculer la rentabilité moyenne des fonds propres pour 2024 et 2025 avec \texttt{in\_options.metrics}, et le produit par salarié de 2025
pour un effectif de 2\,500.
\end{exercise}

\begin{exercise}[$\star\star\star$]\label{exo:in:the-options-market-making-houses:8}
\emph{Trouver la faille.} «~Les maisons d'options gagnent environ 2~millions d'euros par salarié~: les résultats publiés d'une maison le prouvent
pour l'ensemble du secteur.~»
\end{exercise}

\section{Problème~: trois parquets, un écran}

\begin{problem}[{Problème du week-end --- trois parquets, un écran}]\label{pb:in:the-options-market-making-houses:1}
Un étudiant de master qui hésite entre une maison d'options et une société de trading delta-one veut comprendre l'origine des maisons,
leur activité et l'économie d'une maison.

\medskip\noindent\textbf{Partie I --- Origines.}
\begin{enumerate}
\item Définir une \omterm{def:in:the-options-market-making-houses:house}{maison de market making d'options}, et dire en quoi elle diffère du teneur de marché d'options du livre~1.
\item Nommer les trois parquets de la \cref{fig:in:the-options-market-making-houses:lineage} et la société rattachée à chacun.
\item Compter les décennies de fondation des six sociétés.
\item Définir un \omterm{def:in:the-options-market-making-houses:spinout}{essaimage} et dire pourquoi le graphe n'en contient aucun.
\item Que fait encore le parquet de la bourse d'options de Chicago, selon son exploitant~?
\end{enumerate}

\medskip\noindent\textbf{Partie II --- Le travail.}
\begin{enumerate}[resume]
\item Combien de séries un sous-jacent ayant 20 échéances et 50 strikes compte-t-il~?
\item Donner les quatre différences entre une maison d'options et une société delta-one.
\item Qu'enseigne un programme de formation, selon les sociétés, et qu'ajoute-t-il aux portefeuilles~?
\item Pourquoi le logiciel de cotation et le jugement de supervision et de pricing du trader sont-ils une seule et même compétence~?
\item Que lit un candidat pour connaître les conditions d'un programme~?
\end{enumerate}

\medskip\noindent\textbf{Partie III --- Les chiffres d'une maison.}
\begin{enumerate}[resume]
\item Donner le produit net de trading et le résultat net pour 2023--2025.
\item Calculer la croissance du produit net de trading en 2024 et en 2025.
\item Calculer la marge nette de chaque année.
\item Calculer la rentabilité moyenne des fonds propres en 2024 et en 2025.
\item Calculer le produit net de trading par salarié en 2024, et sa borne supérieure en 2025.
\end{enumerate}

\medskip\noindent\textbf{Partie IV --- Le verdict.}
\begin{enumerate}[resume]
\item Énoncer le \emph{résultat nommé}~: le produit net de trading, la marge et la rentabilité moyenne des fonds propres de la maison en 2025, ainsi que
la part des six sociétés rattachées à un parquet selon leur propre récit.
\item Pourquoi «~plus de 2\,000~» est-il une borne supérieure du produit par salarié~?
\item Pourquoi les chiffres publiés d'une seule maison ne peuvent-ils pas valoir pour l'ensemble du secteur~?
\item Quelle source publique renseignerait le mieux l'étudiant sur une maison qui ne publie rien~?
\item En deux phrases, que doit retenir l'étudiant du chapitre pour son choix~?
\end{enumerate}
\end{problem}

\section{Questions d'entretien}

\begin{interviewq}[$\star$ \iqroles{trader}]\label{iq:in:the-options-market-making-houses:1}
Pourquoi un teneur de marché d'options garde-t-il du risque la nuit alors qu'un teneur de marché d'actions cherche à clore sa journée à plat~?
\end{interviewq}

\begin{interviewq}[$\star$ \iqroles{trader, chercheur}]\label{iq:in:the-options-market-making-houses:2}
Un seul titre a des options à 12 échéances et 30 strikes. Combien de cotations un teneur de marché maintient-il
s'il cote les deux côtés de chaque série~?
\end{interviewq}

\begin{interviewq}[$\star\star$ \iqroles{trader}]\label{iq:in:the-options-market-making-houses:3}
La désignation comme teneur de marché par une bourse vous oblige à coter un large éventail de séries. Pourquoi une société accepterait-elle
cette obligation, et que coûte-t-elle~?
\end{interviewq}

\begin{interviewq}[$\star\star$ \iqroles{chercheur, risque}]\label{iq:in:the-options-market-making-houses:4}
Le résultat net d'une société a augmenté de 29~\% et ses fonds propres de 12~\%. Que s'est-il passé pour sa rentabilité des fonds propres, à peu près, et que
vérifieriez-vous avant de conclure que l'activité s'est améliorée~?
\end{interviewq}

\begin{interviewq}[$\star\star$ \iqroles{trader}]\label{iq:in:the-options-market-making-houses:5}
Que vous attendriez-vous à apprendre pendant les six premiers mois dans une maison d'options qu'un manuel n'enseigne pas~?
\end{interviewq}

\begin{interviewq}[$\star\star\star$ \iqroles{trader, risque}]\label{iq:in:the-options-market-making-houses:6}
Votre portefeuille d'options est long en gamma et court en vega après une journée chargée. Décrivez ce que vous feriez avant la clôture, et pourquoi.
\end{interviewq}
